\chapter{Inferência Clássica}
\label{cap:inferencia}
% ── índice ──────────────────────────────────────────────────────
\index{teste t|textbf}
\index{ttest@\texttt{ttest()}|textbf}
\index{ANOVA|textbf}
\index{anova@\texttt{anova()}|textbf}
\index{Kruskal-Wallis, teste de|textbf}
\index{kruskal_wallis@\texttt{kruskal\_wallis()}|textbf}
\index{Welch, teste de|textbf}
\index{não-paramétrico, teste}

% ─────────────────────────────────────────────────────────────────────────────
\section{Testes de hipóteses paramétricos e não-paramétricos}

Os testes apresentados neste capítulo respondem à pergunta: \emph{os dados
são consistentes com uma hipótese sobre a distribuição populacional?}

\subsection*{Teste t de Student}

Para uma amostra de tamanho $n$ com variância desconhecida, a estatística:
\[
  t = \frac{\bar{x} - \mu_0}{s / \sqrt{n}} \sim t_{n-1}
\]
testa $H_0: \mu = \mu_0$. Para duas amostras independentes, o teste de Welch
usa graus de liberdade aproximados pelo método de Satterthwaite, sem assumir
homoscedasticidade.

\subsection*{ANOVA de um fator}

Generaliza o $t$-test para $K > 2$ grupos. Decompõe a variação total:
\[
  \text{SS}_{\text{total}} = \text{SS}_{\text{entre}} + \text{SS}_{\text{dentro}}
\]
\[
  F = \frac{\text{SS}_{\text{entre}} / (K-1)}{\text{SS}_{\text{dentro}} / (N-K)}
  \sim F_{K-1,\, N-K}
\]
$H_0$: todas as médias grupais são iguais.

\subsection*{Kruskal-Wallis}

Alternativa não-paramétrica à ANOVA de um fator. Substitui observações por
\emph{ranks} e compara as somas de ranks entre grupos. Estatística:
\[
  H = \frac{12}{N(N+1)} \sum_{k=1}^K \frac{R_k^2}{n_k} - 3(N+1) \sim \chi^2_{K-1}
\]
válida mesmo sem normalidade, desde que as distribuições tenham forma similar.

% ─────────────────────────────────────────────────────────────────────────────
\section{Verificação por DGP}

DGP com 3 grupos de tamanhos distintos; o valor esperado cresce linearmente
com o grupo. $N=300$, $\sigma = 1$.

\begin{lstlisting}[caption={Testes de hipóteses}]
seed(42)
...
generate df g   = (_n <= 150) * 1 + (_n > 150) * 2
generate df yg  = 2.0 + e + (g - 1) * 0.5  // g=1: mu=2.0, g=2: mu=2.5

// Teste t de uma amostra: H0: mu = 2.0
ttest(df, y, mu=2.0)

// Teste t de dois grupos (Welch): H0: mu(g1) = mu(g2)
ttest(df, yg, by=g)

// Pareado: antes vs depois
ttest(df, y_antes, y_depois, paired=true)

// ANOVA: H0: mu1 = mu2 = mu3
anova(df, y3, by=k)

// Kruskal-Wallis (não-paramétrico)
kruskal_wallis(df, y3, by=k)
\end{lstlisting}

\subsection*{Teste t — uma amostra}

\begin{lstlisting}[language={}, basicstyle=\ttfamily\small, frame=none,
  numbers=none, backgroundcolor=\color{codbg}]
One-sample t-test: y   H0: mean = 2.0
t = -0.6946   df = 299   p = 0.4879
\end{lstlisting}

Não rejeita $H_0$ — média amostral 1{,}98 é compatível com 2{,}0.

\subsection*{Teste t — dois grupos (Welch)}

\begin{lstlisting}[language={}, basicstyle=\ttfamily\small, frame=none,
  numbers=none, backgroundcolor=\color{codbg}]
Two-sample t-test (Welch): yg by g
  grupo 1: n=150  mean=1.971  SE=0.040
  grupo 2: n=150  mean=2.489  SE=0.041
Welch's t = -9.0552   df = 297.38   p = 0.0000
\end{lstlisting}

Rejeita $H_0$: diferença de 0{,}52 ≈ diferença verdadeira de 0{,}50.

\subsection*{ANOVA}

\begin{lstlisting}[language={}, basicstyle=\ttfamily\small, frame=none,
  numbers=none, backgroundcolor=\color{codbg}]
=========================== One-Way ANOVA ============================
Source           SS     df      MS         F       P>F
----------------------------------------------------------------------
Between       51.195     2   25.597   104.037   0.0000
Within        73.074   297    0.246
Total        124.269   299
\end{lstlisting}

$F = 104{,}0$, $p < 0{,}001$: rejeita $H_0$ com segurança.

\subsection*{Kruskal-Wallis}

\begin{lstlisting}[language={}, basicstyle=\ttfamily\small, frame=none,
  numbers=none, backgroundcolor=\color{codbg}]
H = 120.38   df = 2   p = 0.0000
\end{lstlisting}

Confirma a ANOVA sem assumir normalidade.

% ─────────────────────────────────────────────────────────────────────────────
\section{Sintaxe}

\begin{lstlisting}
ttest(df, var, mu=0.0)                  // uma amostra
ttest(df, var, by=grupo)               // dois grupos (Welch)
ttest(df, var1, var2, paired=true)     // pareado
anova(df, var, by=grupo)               // ANOVA de um fator
kruskal_wallis(df, var, by=grupo)      // Kruskal-Wallis
\end{lstlisting}

\begin{nota}
  O teste de Welch não assume homocedasticidade e é preferível ao $t$-test
  clássico de Student em quase todos os casos práticos. A perda de potência
  é negligenciável quando as variâncias são iguais; o ganho quando diferem
  é substancial.
\end{nota}
